% compile this file with xelatex
\documentclass[12pt]{article}
\usepackage{graphicx}
\usepackage[margin=2cm, a4paper]{geometry}
\usepackage{setspace}
\usepackage{pdfpages}
\usepackage{float}
\usepackage{amsmath}
\usepackage{fancyvrb}
\usepackage{amssymb}
\usepackage{minted}
\usepackage{enumitem}
\usepackage[colorlinks,linkcolor=blue]{hyperref}


\renewcommand{\contentsname}{Contents}
\renewcommand{\figurename}{Figure}
\renewcommand{\tablename}{Table}
\hypersetup{
    colorlinks=true,
    linkcolor=black,
    filecolor=magenta,      
    urlcolor=blue,
}
\newcommand{\mytitle}{Numerical Method - Homework 1}
\newcommand{\myauthor}{B13902022 Yu-Chi,Lai}

\usepackage{fancyhdr}
\pagestyle{fancy}
\fancyhead{}
\fancyhead[L]{\mytitle}
\fancyhead[R]{\myauthor}

\title{\mytitle}
\author{\textbf{\myauthor}}
\date{}

\begin{document}

\maketitle

\section{}


\newpage
\section{}
I use the following code to reproduce the example ($f_i=1,A=-64,B=0$) given in the paper, and I got $0$ using the first method (\texttt{calc1()}, $1-p_i=1-\frac{1}{1+e^{Af_i+B}}$), while getting 1.60381e-28 when using the second method (\texttt{calc2()}, $1-p_i=\frac{e^{Af_i+B}}{1+e^{Af_i+B}}$)

The second result is more precise, this is mainly because, when $(f_i,A,B)=(1,-64,0)$, $\mathrm{exp}(Af_i+B)\approx 0$, it successfully avoid the subtraction between two close numbers in \texttt{calc1} (In this case, $1$ and $p_i$), which leads to catastrophic cancellation.
\begin{minted}[frame=lines,framesep=2mm,baselinestretch=1.2,linenos,breaklines]{cpp}
#include <bits/stdc++.h>
using namespace std;

double calc1(double f, double A, double B) {
    return 1.0 - 1.0 / (exp(A*f+B)+1.0);
}

double calc2(double f, double A, double B) {
    return exp(A*f+B) / (1.0+exp(A*f+B));
}

int main() {
    cout << calc1(1, -64, 0) << endl;
    // output: 0
    cout << calc2(1, -64, 0) << endl;
    // output: 1.60381e-28
} 
\end{minted}

\end{document}
